\documentclass[11pt,a4paper]{article}

% ---------- Codificación y tipografía ----------
\usepackage[utf8]{inputenc}
\usepackage[T1]{fontenc}
\usepackage{lmodern}
\usepackage[spanish,es-nodecimaldot]{babel}
\usepackage{microtype}
\usepackage{amsmath,amssymb}
\usepackage{geometry}
\geometry{top=2.3cm,bottom=2.3cm,left=2.2cm,right=2.2cm}
\usepackage{enumitem}
\usepackage{booktabs}
\usepackage{longtable}
\usepackage{array}
\usepackage{fancyhdr}
\usepackage{titlesec}
\usepackage{setspace}
\usepackage{multicol}
\usepackage{needspace}

\setstretch{1.08}
\setlength{\parindent}{0pt}
\setlength{\parskip}{4pt}

% ---------- Encabezado y pie ----------
\pagestyle{fancy}
\fancyhf{}
\lhead{\small Banco de Preguntas PEU --- SUNASS}
\rhead{\small Módulo IV: Equilibrio General y Bienestar}
\cfoot{\small \thepage}
\renewcommand{\headrulewidth}{0.4pt}
\renewcommand{\footrulewidth}{0pt}

% ---------- Títulos ----------
\titleformat{\section}{\large\bfseries}{\thesection}{0.8em}{}
\titlespacing*{\section}{0pt}{14pt}{6pt}
\titleformat{\subsection}{\normalsize\bfseries}{\thesubsection}{0.6em}{}
\titlespacing*{\subsection}{0pt}{10pt}{4pt}

% ---------- Comandos del examen ----------
\newcounter{qnum}
\newcommand{\q}[1]{%
  \needspace{8\baselineskip}%
  \par\medskip\refstepcounter{qnum}%
  \noindent\textbf{\theqnum.}\ #1\par\nopagebreak[4]
}
\newenvironment{opts}
  {\begin{enumerate}[label=\Alph*),leftmargin=2.0em,itemsep=1pt,topsep=2pt,parsep=0pt,partopsep=0pt]}
  {\end{enumerate}}
\makeatletter
\@itempenalty=10000
\makeatother
\newcommand{\sol}[3]{\par\noindent\textbf{#1. #2.}\ #3\par\smallskip}

\begin{document}

% ============================================================
% PORTADA
% ============================================================
\begin{titlepage}
\thispagestyle{empty}
\centering
\vspace*{1.2cm}

{\Large\bfseries BANCO DE PREGUNTAS}\\[0.35cm]
{\large EXAMEN DE ADMISIÓN}\\[0.15cm]
{\large PROGRAMA DE EXTENSIÓN UNIVERSITARIA (PEU)}\\[0.15cm]
{\large SUPERINTENDENCIA NACIONAL DE SERVICIOS DE SANEAMIENTO}\\[0.05cm]
{\large SUNASS}\\[1.0cm]

\rule{\linewidth}{0.9pt}\\[0.55cm]
{\LARGE\bfseries MÓDULO IV: EQUILIBRIO GENERAL Y BIENESTAR}\\[0.35cm]
{\large Especialidad: Economía}\\[0.55cm]
\rule{\linewidth}{0.9pt}\\[1.0cm]

\begin{tabular}{ll}
\toprule
Número de preguntas: & 90 \\
Bloque A: & 30 preguntas \\
Bloque B: & 35 preguntas \\
Bloque C: & 25 preguntas \\
Duración sugerida: & 150 minutos \\
Nivel: & Microeconomía intermedia--avanzada \\
Formato: & Opción múltiple (A, B, C, D) \\
\bottomrule
\end{tabular}

\vfill
{\small Material de preparación académica. Uso interno.}\\[0.15cm]
{\small Edición 2026}
\end{titlepage}

% ============================================================
% INSTRUCCIONES
% ============================================================
\section*{Instrucciones generales}

\begin{enumerate}[label=\arabic*.,leftmargin=1.6em,itemsep=3pt]
  \item El presente banco contiene \textbf{90 preguntas de opción múltiple}. Cada pregunta tiene una sola respuesta correcta.
  \item No se descuenta por respuesta incorrecta; se recomienda responder todas las preguntas.
  \item El tiempo sugerido de resolución es de \textbf{150 minutos}. Use calculadora no programable.
  \item Las preguntas están organizadas en tres bloques: fundamentos teóricos, aplicaciones cuantitativas y casos regulatorios del sector saneamiento.
  \item El \textbf{solucionario completo} y la \textbf{tabla de claves rápidas} se encuentran al final. Se recomienda resolver el examen antes de consultarlos.
\end{enumerate}

\section*{Estructura del banco}

\begin{center}
\renewcommand{\arraystretch}{1.25}
\begin{tabular}{@{}p{1.4cm}>{\raggedright\arraybackslash}p{7.6cm}cc@{}}
\toprule
\textbf{Bloque} & \textbf{Contenido} & \textbf{Preguntas} & \textbf{Peso} \\
\midrule
A & Fundamentos teóricos y conceptuales: equilibrio general, ley de Walras, eficiencia de Pareto, caja de Edgeworth, curva de contrato, núcleo y teoremas del bienestar. & 1--30 & 33.3\% \\
B & Aplicaciones cuantitativas: economías de intercambio, precios relativos, curvas de contrato, frontera de utilidad, bienestar social, externalidades y bienes públicos. & 31--65 & 38.9\% \\
C & Casos con implicancias regulatorias: asignación de agua, derechos de agua, externalidades, monopolio natural, transferencias lump-sum y evaluación de bienestar. & 66--90 & 27.8\% \\
\bottomrule
\end{tabular}
\end{center}

\section*{Competencias evaluadas}

\noindent El banco evalúa la capacidad de interpretar el equilibrio general y la interdependencia de los mercados, analizar la eficiencia de Pareto mediante la caja de Edgeworth, derivar precios y asignaciones de equilibrio en economías de intercambio, aplicar los teoremas del bienestar y evaluar políticas de asignación de recursos y redistribución en el sector saneamiento.

\section*{Notación principal}

{\small
\renewcommand{\arraystretch}{1.0}
\begin{center}
\begin{tabular}{@{}ll@{\hspace{1.0cm}}ll@{}}
\toprule
$U_A, U_B$ & Utilidad de los consumidores A y B & $x_A, y_A$ & Cantidades del consumidor A \\
$X, Y$ & Dotaciones totales de los bienes & $p_x, p_y$ & Precios de los bienes \\
$MRS$ & Tasa marginal de sustitución & $MRT$ & Tasa marginal de transformación \\
$z_i(p)$ & Exceso de demanda del bien $i$ & $W$ & Función de bienestar social \\
$FPP$ & Frontera de posibilidades de producción & $FPU$ & Frontera de posibilidades de utilidad \\
$EC$ & Excedente del consumidor & $EP$ & Excedente del productor \\
\bottomrule
\end{tabular}
\end{center}
}

\newpage
\setcounter{qnum}{0}

\section*{Bloque A. Fundamentos teóricos y conceptuales}

\noindent\textit{Preguntas 1 a 30. Evalúan la comprensión del equilibrio general, la eficiencia de Pareto, la caja de Edgeworth, la curva de contrato, el núcleo y los teoremas fundamentales del bienestar.}

\q{El equilibrio general se define como:}
\begin{opts}
  \item el equilibrio de un solo mercado, sin interacción con los demás.
  \item una situación en la que solo los consumidores maximizan su utilidad.
  \item una asignación fijada por el regulador.
  \item una situación en la que todos los mercados se vacían simultáneamente y los precios se ajustan.
\end{opts}

\q{Un equilibrio walrasiano es:}
\begin{opts}
  \item cualquier asignación Pareto eficiente.
  \item un vector de precios tal que la oferta iguala la demanda en todos los mercados.
  \item un equilibrio que solo considera el lado de la producción.
  \item un precio fijado por el Estado.
\end{opts}

\q{En equilibrio general, solo importan los precios relativos porque:}
\begin{opts}
  \item el dinero es neutral en el corto plazo.
  \item los precios absolutos siempre son cero.
  \item las demandas y ofertas dependen de los cocientes de precios y se puede elegir un numerario.
  \item el Estado fija todos los precios.
\end{opts}

\q{La ley de Walras establece que:}
\begin{opts}
  \item la suma del valor de los excesos de demanda es cero a los precios vigentes.
  \item todos los mercados tienen exceso de demanda positivo.
  \item los precios son iguales en todos los mercados.
  \item la oferta siempre es igual a la demanda en cada mercado.
\end{opts}

\q{La función de exceso de demanda es homogénea de grado:}
\begin{opts}
  \item 1 en precios.
  \item 2 en precios.
  \item $-1$ en precios.
  \item 0 en precios.
\end{opts}

\q{La existencia de un equilibrio competitivo requiere, entre otros supuestos:}
\begin{opts}
  \item externalidades negativas.
  \item preferencias convexas y funciones de demanda continuas.
  \item poder de mercado de una empresa.
  \item precios rígidos fijados por el gobierno.
\end{opts}

\q{La principal diferencia entre el equilibrio parcial y el equilibrio general es que:}
\begin{opts}
  \item el parcial ignora el mercado del bien analizado.
  \item el general estudia un solo mercado.
  \item el general incorpora la interdependencia entre todos los mercados.
  \item el parcial supone competencia imperfecta.
\end{opts}

\q{Una asignación es Pareto eficiente si:}
\begin{opts}
  \item todos los consumidores reciben la misma cantidad.
  \item se maximiza la producción total.
  \item no existe comercio entre los agentes.
  \item no es posible mejorar a alguien sin empeorar a otro.
\end{opts}

\q{Una mejora de Pareto es:}
\begin{opts}
  \item un cambio en el que alguien mejora y otro empeora.
  \item un cambio en el que todos empeoran.
  \item una decisión adoptada por el regulador.
  \item un cambio en el que al menos uno mejora y nadie empeora.
\end{opts}

\q{El Primer Teorema del Bienestar establece que:}
\begin{opts}
  \item todo equilibrio competitivo es Pareto eficiente bajo ciertos supuestos.
  \item toda asignación eficiente es equitativa.
  \item los mercados siempre fallan.
  \item el gobierno debe fijar los precios.
\end{opts}

\q{El Primer Teorema del Bienestar puede fallar cuando:}
\begin{opts}
  \item hay competencia perfecta.
  \item existen externalidades, bienes públicos, poder de mercado o información asimétrica.
  \item los mercados son completos.
  \item todos los agentes son precio-aceptantes.
\end{opts}

\q{El Segundo Teorema del Bienestar establece que:}
\begin{opts}
  \item todo equilibrio competitivo es equitativo.
  \item no se puede alcanzar la eficiencia sin monopolio.
  \item toda asignación Pareto eficiente puede sostenerse como equilibrio competitivo con transferencias lump-sum, bajo ciertos supuestos.
  \item la eficiencia exige precios fijados por el gobierno.
\end{opts}

\q{El Segundo Teorema del Bienestar utiliza transferencias lump-sum porque:}
\begin{opts}
  \item no distorsionan las decisiones marginales y solo redistribuyen ingreso.
  \item siempre aumentan la eficiencia.
  \item las pagan las empresas.
  \item requieren precios máximos.
\end{opts}

\q{La caja de Edgeworth representa:}
\begin{opts}
  \item la frontera de posibilidades de producción.
  \item las asignaciones de dos bienes entre dos consumidores para dotaciones totales dadas.
  \item solo los precios de equilibrio.
  \item el presupuesto del regulador.
\end{opts}

\q{El punto de dotación en la caja de Edgeworth corresponde a:}
\begin{opts}
  \item la curva de contrato.
  \item la asignación inicial de recursos.
  \item el equilibrio de mercado.
  \item el óptimo social.
\end{opts}

\q{La curva de contrato es:}
\begin{opts}
  \item el conjunto de asignaciones Pareto eficientes.
  \item la recta presupuestaria.
  \item la curva de demanda del mercado.
  \item la función de producción.
\end{opts}

\q{El núcleo o core de una economía es:}
\begin{opts}
  \item el conjunto de asignaciones que ninguna coalición puede mejorar.
  \item el conjunto de todas las asignaciones factibles.
  \item el conjunto de las asignaciones igualitarias.
  \item el conjunto de asignaciones fijadas por el gobierno.
\end{opts}

\q{En la caja de Edgeworth, el área en forma de lente representa:}
\begin{opts}
  \item asignaciones ineficientes.
  \item únicamente la dotación inicial.
  \item la curva de contrato.
  \item asignaciones que mejoran a ambos consumidores respecto de la dotación.
\end{opts}

\q{En un óptimo de Pareto interior de intercambio puro se cumple:}
\begin{opts}
  \item $MRS_A>MRS_B$.
  \item $MRS_A<MRS_B$.
  \item $MRS=0$.
  \item $MRS_A=MRS_B$.
\end{opts}

\q{En una economía con producción, la eficiencia requiere que:}
\begin{opts}
  \item $MRS=0$.
  \item $MRS=MRT$.
  \item $MRS>MRT$.
  \item $P=CMe$.
\end{opts}

\q{La frontera de posibilidades de utilidad muestra:}
\begin{opts}
  \item los costos de producción.
  \item los precios de equilibrio.
  \item la máxima utilidad de un consumidor dada la utilidad del otro.
  \item la dotación inicial de recursos.
\end{opts}

\q{Una función de bienestar social:}
\begin{opts}
  \item describe solo las preferencias individuales.
  \item agrega las utilidades individuales en un ordenamiento social.
  \item siempre es Pareto eficiente.
  \item ignora por completo la equidad.
\end{opts}

\q{Una asignación Pareto eficiente:}
\begin{opts}
  \item no necesariamente es equitativa.
  \item siempre es igualitaria.
  \item maximiza la igualdad.
  \item exige que las utilidades sean iguales.
\end{opts}

\q{Una transferencia lump-sum:}
\begin{opts}
  \item distorsiona los precios relativos.
  \item es un impuesto al consumo.
  \item redistribuye ingreso sin crear distorsiones marginales.
  \item es un precio máximo.
\end{opts}

\q{El criterio de Kaldor-Hicks:}
\begin{opts}
  \item exige que nadie empeore.
  \item es idéntico a una mejora de Pareto.
  \item ignora la eficiencia.
  \item considera una mejora potencial de Pareto si los ganadores pueden compensar a los perdedores.
\end{opts}

\q{Una externalidad negativa genera:}
\begin{opts}
  \item ineficiencia asignativa y una falla del Primer Teorema del Bienestar.
  \item eficiencia automática.
  \item equidad.
  \item precios competitivos.
\end{opts}

\q{Un bien público puro se caracteriza por:}
\begin{opts}
  \item no rivalidad y no exclusión.
  \item rivalidad y exclusión.
  \item ser producido por un monopolio.
  \item tener un precio positivo.
\end{opts}

\q{Un monopolio genera pérdida de bienestar porque:}
\begin{opts}
  \item fija $P=CMg$.
  \item produce más que la competencia.
  \item fija $P>CMg$ y produce menos que la cantidad competitiva.
  \item no tiene costos.
\end{opts}

\q{El Primer Teorema del Bienestar requiere, entre otros supuestos:}
\begin{opts}
  \item mercados completos, competencia perfecta, información perfecta y ausencia de externalidades.
  \item monopolio.
  \item externalidades.
  \item precios rígidos.
\end{opts}

\q{Si los mercados son incompletos o hay información asimétrica:}
\begin{opts}
  \item el Primer Teorema se cumple siempre.
  \item la eficiencia aumenta.
  \item el equilibrio competitivo puede no ser Pareto eficiente.
  \item no hay intercambio.
\end{opts}

\section*{Bloque B. Aplicaciones cuantitativas y optimización}

\noindent\textit{Preguntas 31 a 65. Requieren resolver economías de intercambio, calcular precios relativos, identificar la curva de contrato, evaluar mejoras de Pareto y aplicar funciones de bienestar social.}

\noindent\textit{Considere una economía de intercambio puro con dos consumidores A y B, dos bienes $x$ e $y$, dotaciones totales $X=10$ y $Y=10$. Las preferencias son $U_A=x_Ay_A$ y $U_B=x_By_B$. La dotación inicial es A$=(8,2)$ y B$=(2,8)$. Las preguntas 31 a 44 se refieren a esta economía.}

\q{La cantidad total del bien $x$ disponible en la caja de Edgeworth es:}
\begin{opts}
    \item 8
    \item 9
    \item 10
    \item 20
\end{opts}

\q{La cantidad total del bien $y$ disponible en la caja de Edgeworth es:}
\begin{opts}
    \item 8
    \item 9
    \item 12
    \item 10
\end{opts}

\q{La dotación inicial del consumidor A es:}
\begin{opts}
    \item $(2,8)$
    \item $(8,2)$
    \item $(5,5)$
    \item $(10,10)$
\end{opts}

\q{En la dotación inicial, la tasa marginal de sustitución del consumidor A es:}
\begin{opts}
    \item 4.00
    \item 2.00
    \item 1.00
    \item 0.25
\end{opts}

\q{En la dotación inicial, la tasa marginal de sustitución del consumidor B es:}
\begin{opts}
    \item 0.25
    \item 1
    \item 2
    \item 4
\end{opts}

\q{¿La dotación inicial es Pareto eficiente?}
\begin{opts}
    \item Sí, porque ambas tasas marginales de sustitución ya son iguales.
    \item No, porque difieren las RMS y hay ganancias de intercambio.
    \item Sí, porque las dotaciones totales de ambos bienes son fijas.
    \item No, porque los precios relativos iniciales son iguales a cero.
\end{opts}

\q{La curva de contrato de esta economía está dada por:}
\begin{opts}
    \item $y_A=2x_A$.
    \item $y_A=10-x_A$.
    \item $y_A=0$.
    \item $y_A=x_A$.
\end{opts}

\q{El precio relativo de equilibrio $p_x/p_y$ es:}
\begin{opts}
    \item 0.5
    \item 1
    \item 2
    \item 4
\end{opts}

\q{En el equilibrio competitivo, la cantidad del bien $x$ que consume A es:}
\begin{opts}
    \item 4
    \item 4.5
    \item 5
    \item 6
\end{opts}

\q{En el equilibrio competitivo, la cantidad del bien $y$ que consume A es:}
\begin{opts}
    \item 4
    \item 4.5
    \item 5
    \item 6
\end{opts}

\q{La utilidad de A en la dotación inicial es 16 y en el equilibrio es 25. La ganancia de utilidad de A es:}
\begin{opts}
    \item 0
    \item 5
    \item 7
    \item 9
\end{opts}

\q{La ganancia de utilidad del consumidor B al pasar de la dotación al equilibrio es:}
\begin{opts}
    \item 0
    \item 5
    \item 9
    \item 16
\end{opts}

\q{El equilibrio competitivo es Pareto eficiente porque:}
\begin{opts}
    \item A consume más que B.
    \item B pierde con el intercambio.
    \item se cumple $MRS_A=MRS_B=1$.
    \item los precios absolutos son iguales.
\end{opts}

\q{Si se desea alcanzar la asignación eficiente A$=(6,6)$ y B$=(4,4)$ con $p_x=p_y=1$, la transferencia lump-sum de B hacia A debe ser:}
\begin{opts}
    \item 2
    \item 1
    \item 4
    \item 0
\end{opts}

\q{Considere la misma economía, pero con dotación inicial A$=(9,1)$ y B$=(1,9)$. Si $p_x/p_y=1$, la cantidad de $x$ que consume A en equilibrio es:}
\begin{opts}
    \item 4
    \item 6
    \item 8
    \item 5
\end{opts}

\q{Con la dotación A$=(9,1)$ y B$=(1,9)$, si se transfiere 2 de B hacia A, la cantidad de $x$ que consume A es:}
\begin{opts}
    \item 4
    \item 5
    \item 6
    \item 7
\end{opts}

\q{La asignación A$=(6,6)$ y B$=(4,4)$ es Pareto eficiente porque:}
\begin{opts}
    \item las RMS de ambos consumidores son iguales.
    \item el consumidor B recibe menor cantidad que A.
    \item la dotación agregada permanece fija en el tiempo.
    \item los precios relativos de equilibrio valen cero.
\end{opts}

\q{El Segundo Teorema del Bienestar implica que:}
\begin{opts}
    \item toda asignación eficiente se sostiene con precios y transferencias lump-sum.
    \item todo equilibrio competitivo garantiza automáticamente equidad distributiva.
    \item la eficiencia asignativa exige necesariamente propiedad estatal de empresas.
    \item las transferencias lump-sum distorsionan los precios relativos de mercado.
\end{opts}

\q{En una caja de Edgeworth con $X=20$ e $Y=10$, si A tiene $(12,4)$, la cantidad de $x$ que tiene B es:}
\begin{opts}
    \item 4
    \item 6
    \item 10
    \item 8
\end{opts}

\q{En la misma caja, si A tiene $(12,4)$, la cantidad de $y$ que tiene B es:}
\begin{opts}
    \item 4
    \item 6
    \item 8
    \item 10
\end{opts}

\q{Si en una asignación interior con preferencias convexas se cumple $MRS_A=MRS_B=2$, entonces:}
\begin{opts}
    \item la asignación es Pareto eficiente.
    \item la asignación es ineficiente.
    \item solo es eficiente si los precios son iguales.
    \item solo es eficiente si las dotaciones son iguales.
\end{opts}

\q{Si $MRS_A=3$ y $MRS_B=1$, entonces:}
\begin{opts}
    \item no existen ganancias posibles derivadas del intercambio.
    \item la asignación actual ya es estrictamente Pareto eficiente.
    \item solo una expropiación directa mejora el bienestar conjunto.
    \item hay ganancias de intercambio al diferir las valoraciones.
\end{opts}

\q{Con dos bienes, $p_x=p_y=1$ y un exceso de demanda $z_x=4$, la ley de Walras implica que $z_y$ es:}
\begin{opts}
    \item $-4$
    \item $+0$
    \item $+4$
    \item $+8$
\end{opts}

\q{Si $p_x=p_y=1$ y $z_x=12$, entonces $z_y$ es:}
\begin{opts}
    \item $+12$
    \item $+00$
    \item $-12$
    \item $+24$
\end{opts}

\q{Si $p_x/p_y=2$ y $p_y=1$, entonces $p_x$ es:}
\begin{opts}
    \item 0.5
    \item 1
    \item 2
    \item 4
\end{opts}

\q{Si $p_x=3$ y $p_y=6$, el precio relativo $p_x/p_y$ es:}
\begin{opts}
    \item 2.0
    \item 0.5
    \item 3.0
    \item 6.0
\end{opts}

\q{Si la frontera de posibilidades de utilidad es $U_A+U_B=100$, el máximo valor de $U_A$ es:}
\begin{opts}
    \item 0
    \item 50
    \item 100
    \item 200
\end{opts}

\q{Si $W=\min(U_A,U_B)$ y $U_A+U_B=100$, el máximo de $W$ se alcanza con $U_A$ igual a:}
\begin{opts}
    \item 0
    \item 25
    \item 50
    \item 100
\end{opts}

\q{Si $W=U_A+U_B$ y $U_A+U_B=100$, el máximo de $W$ es:}
\begin{opts}
    \item 50
    \item 0
    \item 200
    \item 100
\end{opts}

\q{Una función de bienestar social que maximiza la utilidad del individuo peor situado corresponde a:}
\begin{opts}
    \item una función utilitarista.
    \item una función rawlsiana.
    \item el criterio de Pareto.
    \item el criterio de Kaldor-Hicks.
\end{opts}

\q{El Primer Teorema del Bienestar requiere:}
\begin{opts}
    \item mercados completos, sin externalidades y precio-aceptantes.
    \item presencia de monopolio natural y barreras legales a la entrada.
    \item existencia de externalidades negativas y bienes públicos puros.
    \item precios máximos uniformes fijados directamente por el Estado.
\end{opts}

\q{Si existe una externalidad negativa, el equilibrio competitivo:}
\begin{opts}
    \item es siempre Pareto eficiente.
    \item no es Pareto eficiente.
    \item es necesariamente equitativo.
    \item tiene exceso de demanda cero.
\end{opts}

\q{Un impuesto pigouviano óptimo debe igualar:}
\begin{opts}
    \item el costo externo marginal.
    \item el costo medio total.
    \item el costo total del bien.
    \item cero.
\end{opts}

\q{La condición de Samuelson para la provisión eficiente de un bien público es:}
\begin{opts}
    \item $MRS_A = MRS_B$
    \item $MRS_i = MRT_i$
    \item $\sum MRS = MRT$
    \item $\sum MRS_i = 0$
\end{opts}

\q{En un monopolio con $P>CMg$, el resultado es:}
\begin{opts}
    \item eficiencia asignativa plena en todos los mercados.
    \item ausencia absoluta de pérdida irrecuperable de bienestar.
    \item pérdida de bienestar y falla del Primer Teorema.
    \item un equilibrio competitivo Pareto eficiente sin costos.
\end{opts}

\section*{Bloque C. Casos con implicancias regulatorias y de política}

\noindent\textit{Preguntas 66 a 90. Aplican el equilibrio general, la eficiencia de Pareto, la caja de Edgeworth y los teoremas del bienestar a la asignación de agua, los derechos de agua, las externalidades, el monopolio natural y la evaluación de políticas tarifarias.}

\noindent\textit{Para las preguntas 66 a 69, considere dos usuarios de agua, A y B, con un total de 100 unidades. Los beneficios marginales son $MBA=100-q_A$ y $MBB=80-q_B$, y los beneficios totales son las áreas bajo estas curvas.}

\q{La asignación eficiente de agua para el usuario A es:}
\begin{opts}
  \item 50
  \item 60
  \item 70
  \item 80
\end{opts}

\q{La asignación eficiente de agua para el usuario B es:}
\begin{opts}
  \item 30
  \item 40
  \item 50
  \item 60
\end{opts}

\q{Si la asignación inicial es A$=70$ y B$=30$, la ganancia total de beneficio al pasar a la asignación eficiente es:}
\begin{opts}
  \item 50
  \item 75
  \item 100
  \item 200
\end{opts}

\q{En una caja de Edgeworth con 100 unidades de agua y 200 unidades del otro bien, si A tiene $(60,80)$, la cantidad del otro bien que tiene B es:}
\begin{opts}
  \item 80
  \item 120
  \item 100
  \item 60
\end{opts}

\q{Si en una asignación de agua $MRS_A=3$ y $MRS_B=1$, entonces:}
\begin{opts}
  \item ambos usuarios pueden ganar con el intercambio.
  \item no existen ganancias de intercambio.
  \item solo pueden ganar si los precios son iguales.
  \item solo el regulador puede mejorar la asignación.
\end{opts}

\q{En un óptimo de Pareto interior de asignación de agua entre dos usuarios se cumple:}
\begin{opts}
  \item $MRS_A>MRS_B$.
  \item $MRS_A=MRS_B$.
  \item $MRS=0$.
  \item $MRS=MRT$ únicamente.
\end{opts}

\q{Si una asignación de agua se encuentra sobre la curva de contrato, entonces:}
\begin{opts}
  \item no es eficiente.
  \item es necesariamente equitativa.
  \item maximiza la producción física.
  \item es Pareto eficiente.
\end{opts}

\q{Si la asignación inicial es A$=80$ y B$=20$, y la asignación eficiente es A$=60$ y B$=40$, la transferencia de agua de A hacia B es:}
\begin{opts}
  \item 20
  \item 10
  \item 30
  \item 40
\end{opts}

\q{La reasignación de 20 unidades desde A hacia B aumenta el beneficio total. Esto implica que:}
\begin{opts}
  \item existen ganancias potenciales de bienestar; con compensación puede ser una mejora de Pareto.
  \item A mejora con seguridad aunque no reciba compensación.
  \item B necesariamente pierde con la reasignación.
  \item no existen ganancias de bienestar.
\end{opts}

\q{Si los derechos de agua están bien definidos y los costos de transacción son cero, el teorema de Coase implica que:}
\begin{opts}
  \item solo el gobierno puede asignar el agua eficientemente.
  \item la asignación eficiente requiere repartir el agua por igual.
  \item no habrá intercambio de derechos.
  \item se puede alcanzar una asignación eficiente con independencia de la asignación inicial de derechos.
\end{opts}

\q{Para que el Primer Teorema del Bienestar se aplique a un mercado de agua se requiere, entre otros supuestos:}
\begin{opts}
  \item existencia de monopolio y barreras a la entrada.
  \item ausencia de externalidades, mercados completos y agentes precio-aceptantes.
  \item externalidades negativas y bienes públicos.
  \item precios máximos fijados por el regulador.
\end{opts}

\q{Si la extracción de agua genera una externalidad negativa, el equilibrio competitivo:}
\begin{opts}
  \item es siempre Pareto eficiente.
  \item es necesariamente equitativo.
  \item no es Pareto eficiente.
  \item tiene exceso de demanda cero.
\end{opts}

\q{Un impuesto pigouviano sobre la extracción de agua debe fijarse igual a:}
\begin{opts}
  \item cero.
  \item el costo medio total.
  \item el costo total del servicio.
  \item el costo externo marginal.
\end{opts}

\q{El Segundo Teorema del Bienestar aplicado al agua implica que:}
\begin{opts}
  \item todo equilibrio competitivo es equitativo.
  \item una asignación eficiente puede sostenerse con precios y transferencias lump-sum, bajo ciertos supuestos.
  \item la eficiencia exige un monopolio estatal.
  \item se debe fijar un precio máximo.
\end{opts}

\q{Para proteger a los hogares vulnerables sin distorsionar las señales de precio, la opción más adecuada es:}
\begin{opts}
  \item un precio máximo general por debajo del equilibrio.
  \item un precio mínimo.
  \item una cuota de consumo.
  \item una transferencia lump-sum focalizada.
\end{opts}

\q{En un monopolio natural de agua, fijar el precio igual al costo marginal genera:}
\begin{opts}
  \item beneficios extraordinarios.
  \item beneficio económico nulo.
  \item ausencia de pérdidas.
  \item pérdidas que requieren un subsidio o una tarifa en dos partes.
\end{opts}

\q{La protección de una cuenca como bien público se caracteriza por:}
\begin{opts}
  \item ser rival y excluible.
  \item ser provista en exceso por el mercado.
  \item no ser rival ni excluible, por lo que el mercado tiende a subproveerla.
  \item ser provista eficientemente por una empresa competitiva.
\end{opts}

\q{La condición de Samuelson para la provisión eficiente de un bien público es:}
\begin{opts}
  \item $\sum MRS=MRT$.
  \item $MRS=0$.
  \item $MRS=MRT$.
  \item $\sum MRS=0$.
\end{opts}

\q{El criterio de Kaldor-Hicks sostiene que una política es deseable si:}
\begin{opts}
  \item nadie empeora.
  \item los ganadores podrían compensar a los perdedores y aun así mejorar.
  \item todos reciben la misma utilidad.
  \item el gobierno fija los precios.
\end{opts}

\q{Si los ganadores compensan efectivamente a los perdedores y nadie queda peor, el resultado es:}
\begin{opts}
  \item una mejora de Pareto.
  \item solo una mejora potencial de Kaldor-Hicks.
  \item una asignación ineficiente.
  \item un precio máximo.
\end{opts}

\q{Una asignación Pareto eficiente:}
\begin{opts}
  \item es siempre equitativa.
  \item puede ser eficiente pero inequitativa, por lo que se requiere un juicio distributivo.
  \item siempre iguala las utilidades.
  \item elimina la necesidad de redistribución.
\end{opts}

\q{Una función de bienestar social rawlsiana se concentra en:}
\begin{opts}
  \item la suma de las utilidades.
  \item el criterio de Kaldor-Hicks.
  \item la utilidad del individuo peor situado.
  \item la eficiencia productiva.
\end{opts}

\q{Una función de bienestar social utilitarista busca maximizar:}
\begin{opts}
  \item la suma de las utilidades individuales.
  \item la utilidad del peor situado.
  \item la desigualdad.
  \item el precio del agua.
\end{opts}

\q{Al evaluar una reforma tarifaria de agua en equilibrio general, deben considerarse:}
\begin{opts}
  \item únicamente el mercado de agua.
  \item únicamente a los consumidores.
  \item únicamente a los productores.
  \item los efectos en otros mercados, como alimentos, energía y trabajo.
\end{opts}

\q{La evaluación de bienestar de una política de agua debe incluir:}
\begin{opts}
  \item solo el excedente del consumidor.
  \item las externalidades y los efectos de equilibrio general, además de los excedentes.
  \item solo el excedente del productor.
  \item solo la recaudación fiscal.
\end{opts}

\newpage
\section*{Tabla de claves rápidas}

\noindent Use esta tabla para una verificación inmediata. El solucionario con la justificación de cada respuesta se presenta en la sección siguiente.

\vspace{0.5cm}
\begin{center}
\renewcommand{\arraystretch}{1.35}
\begin{tabular}{|c|c|c|c|c|c|c|c|c|c|c|}
\hline
\textbf{N.} & \textbf{1} & \textbf{2} & \textbf{3} & \textbf{4} & \textbf{5} & \textbf{6} & \textbf{7} & \textbf{8} & \textbf{9} & \textbf{10} \\
\hline
\textbf{Clave} & D & B & C & A & D & B & C & D & D & A \\
\hline
\textbf{N.} & \textbf{11} & \textbf{12} & \textbf{13} & \textbf{14} & \textbf{15} & \textbf{16} & \textbf{17} & \textbf{18} & \textbf{19} & \textbf{20} \\
\hline
\textbf{Clave} & B & C & A & B & B & A & A & D & D & B \\
\hline
\textbf{N.} & \textbf{21} & \textbf{22} & \textbf{23} & \textbf{24} & \textbf{25} & \textbf{26} & \textbf{27} & \textbf{28} & \textbf{29} & \textbf{30} \\
\hline
\textbf{Clave} & C & B & A & C & D & A & A & C & A & C \\
\hline
\textbf{N.} & \textbf{31} & \textbf{32} & \textbf{33} & \textbf{34} & \textbf{35} & \textbf{36} & \textbf{37} & \textbf{38} & \textbf{39} & \textbf{40} \\
\hline
\textbf{Clave} & C & D & B & D & D & B & D & B & C & C \\
\hline
\textbf{N.} & \textbf{41} & \textbf{42} & \textbf{43} & \textbf{44} & \textbf{45} & \textbf{46} & \textbf{47} & \textbf{48} & \textbf{49} & \textbf{50} \\
\hline
\textbf{Clave} & D & C & C & A & D & C & A & A & D & B \\
\hline
\textbf{N.} & \textbf{51} & \textbf{52} & \textbf{53} & \textbf{54} & \textbf{55} & \textbf{56} & \textbf{57} & \textbf{58} & \textbf{59} & \textbf{60} \\
\hline
\textbf{Clave} & A & D & A & C & C & B & C & C & D & B \\
\hline
\textbf{N.} & \textbf{61} & \textbf{62} & \textbf{63} & \textbf{64} & \textbf{65} & \textbf{66} & \textbf{67} & \textbf{68} & \textbf{69} & \textbf{70} \\
\hline
\textbf{Clave} & A & B & A & C & C & B & B & C & B & A \\
\hline
\textbf{N.} & \textbf{71} & \textbf{72} & \textbf{73} & \textbf{74} & \textbf{75} & \textbf{76} & \textbf{77} & \textbf{78} & \textbf{79} & \textbf{80} \\
\hline
\textbf{Clave} & B & D & A & A & D & B & C & D & B & D \\
\hline
\textbf{N.} & \textbf{81} & \textbf{82} & \textbf{83} & \textbf{84} & \textbf{85} & \textbf{86} & \textbf{87} & \textbf{88} & \textbf{89} & \textbf{90} \\
\hline
\textbf{Clave} & D & C & A & B & A & B & C & A & D & B \\
\hline
\end{tabular}
\end{center}

\vspace{0.4cm}
\noindent\textbf{Distribución de claves:} A = 22; B = 23; C = 23; D = 22.

\newpage
\section*{Solucionario completo}

\noindent Cada respuesta incluye la clave y una justificación breve. En los ejercicios numéricos se indica la ecuación o el cálculo principal.

\subsection*{Bloque A. Fundamentos teóricos y conceptuales}

\sol{1}{D}{El equilibrio general considera simultáneamente todos los mercados de la economía. Los precios se ajustan hasta que la oferta y la demanda coinciden en cada mercado.}

\sol{2}{B}{Un equilibrio walrasiano es un vector de precios tal que todos los mercados se vacían. No es cualquier asignación eficiente ni un precio fijado por el Estado.}

\sol{3}{C}{Las funciones de demanda y oferta dependen de los precios relativos. Además, se puede elegir un bien como numerario, de modo que solo importan $k-1$ precios independientes.}

\sol{4}{A}{La ley de Walras indica que la suma del valor de los excesos de demanda es cero: $\sum_i p_i z_i(p)=0$. Si $n-1$ mercados se vacían, el mercado restante también lo hace.}

\sol{5}{D}{La función de exceso de demanda es homogénea de grado cero en precios: si todos los precios se multiplican por una constante positiva, las cantidades demandadas y ofrecidas no cambian.}

\sol{6}{B}{La existencia de equilibrio requiere continuidad de las funciones de exceso de demanda, convexidad de las preferencias y otros supuestos estándar. El monopolio o los precios rígidos no forman parte de esos supuestos.}

\sol{7}{C}{El equilibrio parcial analiza un mercado aislando los demás; el equilibrio general incorpora las interacciones y los efectos de retroalimentación entre mercados.}

\sol{8}{D}{Una asignación es Pareto eficiente cuando no existe otra asignación factible que mejore a al menos una persona sin empeorar a otra.}

\sol{9}{D}{Una mejora de Pareto exige que al menos un individuo esté mejor y que ningún otro esté peor.}

\sol{10}{A}{El Primer Teorema del Bienestar señala que un equilibrio competitivo es Pareto eficiente bajo supuestos como mercados completos, competencia perfecta y ausencia de externalidades.}

\sol{11}{B}{Las externalidades, los bienes públicos, el poder de mercado y la información asimétrica son fallas de mercado que pueden invalidar el Primer Teorema del Bienestar.}

\sol{12}{C}{El Segundo Teorema del Bienestar establece que, bajo supuestos de convexidad y mercados completos, cualquier asignación Pareto eficiente puede descentralizarse como un equilibrio competitivo con transferencias lump-sum.}

\sol{13}{A}{Las transferencias lump-sum no dependen de las decisiones de los agentes y no crean distorsiones marginales; solo redistribuyen ingreso entre ellos.}

\sol{14}{B}{La caja de Edgeworth representa todas las asignaciones factibles de dos bienes entre dos consumidores, dadas las dotaciones totales.}

\sol{15}{B}{El punto de dotación es la asignación inicial de recursos de la que parten los consumidores antes de intercambiar.}

\sol{16}{A}{La curva de contrato es el conjunto de asignaciones Pareto eficientes. En intercambio puro, corresponde a los puntos donde las curvas de indiferencia son tangentes.}

\sol{17}{A}{El núcleo o core está formado por las asignaciones que ninguna coalición de agentes puede mejorar mediante sus propios recursos.}

\sol{18}{D}{La lente está delimitada por las curvas de indiferencia que pasan por la dotación inicial; todos sus puntos mejoran a ambos consumidores respecto de esa dotación.}

\sol{19}{D}{En un óptimo interior de intercambio puro, las curvas de indiferencia son tangentes y, por tanto, $MRS_A=MRS_B$.}

\sol{20}{B}{Con producción, la eficiencia requiere que la tasa marginal de sustitución entre bienes sea igual a la tasa marginal de transformación: $MRS=MRT$.}

\sol{21}{C}{La frontera de posibilidades de utilidad muestra las combinaciones máximas de utilidad de dos individuos; indica la máxima utilidad de uno dada la utilidad del otro.}

\sol{22}{B}{Una función de bienestar social agrega las utilidades individuales en un ordenamiento social. La elección de la función implica un juicio distributivo.}

\sol{23}{A}{La eficiencia de Pareto se refiere a la imposibilidad de mejorar sin empeorar a otro; no dice nada sobre la igualdad o la justicia de la distribución.}

\sol{24}{C}{Una transferencia lump-sum redistribuye ingreso sin modificar los incentivos marginales ni introducir distorsiones en las decisiones de consumo o producción.}

\sol{25}{D}{El criterio de Kaldor-Hicks acepta una política si los ganadores podrían compensar a los perdedores y aun así mejorar, aunque la compensación no se pague efectivamente.}

\sol{26}{A}{Una externalidad negativa hace que el costo social marginal difiera del costo privado; el equilibrio competitivo no internaliza ese costo y no es Pareto eficiente.}

\sol{27}{A}{Un bien público puro no es rival ni excluible. El mercado tiende a subproveerlo porque los individuos pueden beneficiarse sin pagar.}

\sol{28}{C}{El monopolio fija $P>CMg$ y produce menos que la cantidad competitiva, lo que genera una pérdida de bienestar.}

\sol{29}{A}{El Primer Teorema del Bienestar requiere mercados completos, competencia perfecta, información perfecta, ausencia de externalidades y otros supuestos.}

\sol{30}{C}{Con mercados incompletos o información asimétrica, el equilibrio competitivo puede no ser Pareto eficiente; el Primer Teorema deja de aplicarse automáticamente.}

\subsection*{Bloque B. Aplicaciones cuantitativas y optimización}

\sol{31}{C}{La caja de Edgeworth tiene ancho igual a la dotación total del bien $x$: $X=10$.}

\sol{32}{D}{La altura de la caja de Edgeworth es la dotación total del bien $y$: $Y=10$.}

\sol{33}{B}{La dotación inicial de A es la pareja $(8,2)$. La de B es $(2,8)$, que completa las dotaciones totales.}

\sol{34}{D}{Para $U_A=x_Ay_A$, $MRS_A=MU_x/MU_y=y_A/x_A$. En la dotación: $2/8=0.25$.}

\sol{35}{D}{Para B, $MRS_B=y_B/x_B=8/2=4$.}

\sol{36}{B}{En la dotación, $MRS_A=0.25$ y $MRS_B=4$. Como son distintas, existen ganancias de intercambio y la asignación no es Pareto eficiente.}

\sol{37}{D}{La eficiencia exige $y_A/x_A=y_B/x_B$. Con $x_B=10-x_A$ e $y_B=10-y_A$ se obtiene $y_A=x_A$.}

\sol{38}{B}{Con preferencias Cobb-Douglas idénticas y dotaciones simétricas, el precio relativo de equilibrio es $p_x/p_y=1$.}

\sol{39}{C}{Con $p_x=p_y=1$, el ingreso de A es $8+2=10$. Gasta la mitad en $x$: $x_A=5$.}

\sol{40}{C}{Con el mismo ingreso y precios, A gasta la otra mitad en $y$: $y_A=5$. El equilibrio de A es $(5,5)$.}

\sol{41}{D}{En la dotación, $U_A=8\cdot 2=16$. En el equilibrio, $U_A=5\cdot 5=25$. La ganancia es $25-16=9$.}

\sol{42}{C}{Para B, $U_B=2\cdot 8=16$ en la dotación y $U_B=5\cdot 5=25$ en el equilibrio; la ganancia también es 9.}

\sol{43}{C}{En el equilibrio, A consume $(5,5)$ y B consume $(5,5)$. Entonces $MRS_A=5/5=1$ y $MRS_B=5/5=1$; la asignación es Pareto eficiente.}

\sol{44}{A}{Con $p_x=p_y=1$, la asignación A$=(6,6)$ requiere un ingreso de 12. A tiene un ingreso inicial de 10, por lo que necesita una transferencia de 2 desde B.}

\sol{45}{D}{Con dotación A$=(9,1)$ y B$=(1,9)$, a precios unitarios ambos tienen ingreso 10. A consume $x_A=5$ e $y_A=5$.}

\sol{46}{C}{Con una transferencia de 2 desde B hacia A, el ingreso de A es 12 y el de B es 8. A consume $x_A=6$ e $y_A=6$.}

\sol{47}{A}{La asignación A$=(6,6)$ y B$=(4,4)$ cumple $MRS_A=6/6=1$ y $MRS_B=4/4=1$. Las MRS son iguales, por lo que es Pareto eficiente.}

\sol{48}{A}{El Segundo Teorema del Bienestar permite sostener una asignación Pareto eficiente como equilibrio competitivo mediante transferencias lump-sum, bajo supuestos de convexidad y mercados completos.}

\sol{49}{D}{Si el total de $x$ es 20 y A tiene 12, entonces B tiene $20-12=8$.}

\sol{50}{B}{Si el total de $y$ es 10 y A tiene 4, entonces B tiene $10-4=6$.}

\sol{51}{A}{Con preferencias convexas, en un óptimo interior de intercambio la igualdad de las MRS es la condición de tangencia y, por tanto, de eficiencia de Pareto.}

\sol{52}{D}{Si $MRS_A=3$ y $MRS_B=1$, las valoraciones marginales difieren; existe un intercambio mutuamente beneficioso hasta igualar las MRS.}

\sol{53}{A}{La ley de Walras con dos bienes y precios unitarios implica $z_x+z_y=0$. Si $z_x=4$, entonces $z_y=-4$.}

\sol{54}{C}{Con $p_x=p_y=1$ y $z_x=12$, la ley de Walras implica $z_y=-12$.}

\sol{55}{C}{Si $p_x/p_y=2$ y $p_y=1$, entonces $p_x=2$.}

\sol{56}{B}{El precio relativo es $p_x/p_y=3/6=0.5$.}

\sol{57}{C}{Si $U_A+U_B=100$ y $U_B\geq 0$, el máximo de $U_A$ se alcanza cuando $U_B=0$: $U_A=100$.}

\sol{58}{C}{Para $W=\min(U_A,U_B)$ con $U_A+U_B=100$, el máximo se alcanza en el punto igualitario: $U_A=U_B=50$.}

\sol{59}{D}{Para $W=U_A+U_B$ y $U_A+U_B=100$, cualquier punto de la frontera da $W=100$. La distribución no queda determinada por el criterio utilitarista.}

\sol{60}{B}{La función rawlsiana maximiza la utilidad del individuo peor situado; se concentra en el mínimo de las utilidades.}

\sol{61}{A}{El Primer Teorema del Bienestar requiere mercados completos, ausencia de externalidades y agentes precio-aceptantes, entre otros supuestos.}

\sol{62}{B}{Con una externalidad negativa, el costo social marginal no coincide con el privado; el equilibrio competitivo no es Pareto eficiente.}

\sol{63}{A}{El impuesto pigouviano óptimo iguala el costo externo marginal, de modo que el agente internaliza el daño que genera.}

\sol{64}{C}{La condición de Samuelson para un bien público es que la suma de las tasas marginales de sustitución iguale la tasa marginal de transformación: $\sum MRS=MRT$.}

\sol{65}{C}{En un monopolio, $P>CMg$ y la cantidad producida es menor que la competitiva; esto genera pérdida de bienestar y una falla del Primer Teorema del Bienestar.}

\subsection*{Bloque C. Casos con implicancias regulatorias y de política}

\sol{66}{B}{La eficiencia requiere $MBA=MBB$: $100-q_A=80-q_B$. Con $q_A+q_B=100$, se obtiene $100-q_A=80-(100-q_A)$, de donde $q_A=60$.}

\sol{67}{B}{Con $q_A=60$ y un total de 100, el usuario B recibe $q_B=40$.}

\sol{68}{C}{El beneficio total inicial es $4550+1950=6500$. En la asignación eficiente es $4200+2400=6600$. La ganancia es 100.}

\sol{69}{B}{Si el total del otro bien es 200 y A tiene 80, B tiene $200-80=120$.}

\sol{70}{A}{Con $MRS_A=3$ y $MRS_B=1$, las valoraciones marginales difieren; existe un intercambio que mejora a ambos.}

\sol{71}{B}{En un óptimo interior de asignación de agua, la tangencia de las curvas de indiferencia implica $MRS_A=MRS_B$.}

\sol{72}{D}{La curva de contrato está formada precisamente por las asignaciones Pareto eficientes.}

\sol{73}{A}{La diferencia entre la asignación inicial A$=80$ y la eficiente A$=60$ es 20 unidades, que se transfieren de A hacia B.}

\sol{74}{A}{La reasignación aumenta el beneficio total, pero A podría perder si no recibe compensación. Existen ganancias potenciales que permiten una mejora de Pareto con compensación.}

\sol{75}{D}{Con derechos bien definidos y costos de transacción nulos, el teorema de Coase sostiene que el intercambio conduce a una asignación eficiente, independientemente de la asignación inicial de derechos.}

\sol{76}{B}{El Primer Teorema del Bienestar requiere ausencia de externalidades, mercados completos y agentes precio-aceptantes, entre otros supuestos.}

\sol{77}{C}{La externalidad negativa hace que el costo social marginal supere al privado; el equilibrio competitivo no internaliza el daño y no es Pareto eficiente.}

\sol{78}{D}{El impuesto pigouviano debe igualar el costo externo marginal para corregir la externalidad.}

\sol{79}{B}{El Segundo Teorema del Bienestar permite sostener una asignación eficiente de agua mediante precios y transferencias lump-sum, bajo ciertos supuestos.}

\sol{80}{D}{Una transferencia lump-sum focalizada protege el ingreso de los hogares vulnerables sin distorsionar los precios relativos ni generar escasez.}

\sol{81}{D}{En un monopolio natural, $P=CMg<CMe$ genera pérdidas; se requiere un subsidio, una tarifa en dos partes u otro mecanismo de cobertura.}

\sol{82}{C}{La protección de una cuenca es un bien público si no es rival ni excluible; el mercado tiende a subproveerla porque hay incentivos a no pagar.}

\sol{83}{A}{La condición de Samuelson para un bien público es $\sum MRS=MRT$; la suma de las disposiciones marginales a pagar debe igualar el costo marginal de provisión.}

\sol{84}{B}{El criterio de Kaldor-Hicks evalúa una política como deseable si los ganadores podrían compensar a los perdedores y aun así mejorar, aunque la compensación no se realice.}

\sol{85}{A}{Si la compensación se paga efectivamente y nadie queda peor, el resultado es una mejora de Pareto, no solo una mejora potencial.}

\sol{86}{B}{La eficiencia de Pareto no garantiza equidad. Una asignación puede ser eficiente y, al mismo tiempo, muy desigual; se requiere un juicio distributivo adicional.}

\sol{87}{C}{La función rawlsiana maximiza la utilidad del individuo peor situado, es decir, se concentra en el mínimo de las utilidades.}

\sol{88}{A}{La función utilitarista maximiza la suma de las utilidades individuales, sin atender directamente a su distribución.}

\sol{89}{D}{Una reforma tarifaria puede afectar los precios de otros bienes y los ingresos de los hogares; el análisis de equilibrio general incorpora esos efectos en alimentos, energía y trabajo.}

\sol{90}{B}{La evaluación de bienestar debe incluir externalidades y efectos de equilibrio general, además de los excedentes del consumidor y del productor.}

\subsection*{Bloque D. Ampliación: Teoría y Aplicaciones Prácticas (91 a 180)}

\sol{91}{B}{El equilibrio parcial mantiene ceteris paribus los precios de otros bienes y factores; una reforma amplia altera salarios, rentas y precios relativos en múltiples mercados simultáneamente.}
\sol{92}{C}{Por homogeneidad de grado cero ($z(\lambda p) = z(p)$), solo importan los $k-1$ precios relativos; fijar $p_k = 1$ elimina la indeterminación del nivel absoluto de precios.}
\sol{93}{A}{Bajo no saciedad local, cada consumidor gasta todo su ingreso ($p \cdot x_h = p \cdot \omega_h$). Sumando sobre todos los agentes resulta $p \cdot z(p) \equiv 0$ para cualquier $p > 0$.}
\sol{94}{D}{Como $\sum_{i=1}^k p_i z_i(p) = 0$ y $z_1 = \dots = z_{k-1} = 0$, queda $p_k z_k(p) = 0$. Dado $p_k > 0$, necesariamente $z_k(p) = 0$.}
\sol{95}{B}{La Ley de Walras deriva de las restricciones presupuestarias individuales, por lo que vale para todo vector de precios $p > 0$: un exceso de demanda en un mercado implica exceso de oferta en valor en otros.}
\sol{96}{C}{Por holgura complementaria ($p_j^* z_j(p^*) = 0$ con $z_j \le 0$ y $p_j^* \ge 0$), cualquier bien con exceso de oferta estricto en equilibrio tiene precio nulo ($p_j^* = 0$).}
\sol{97}{A}{En el tâtonnement, $\dot{p}_i = \alpha_i z_i(p)$ y no se ejecutan transacciones hasta que el subastador anuncia el vector de precios que vacía todos los mercados.}
\sol{98}{D}{Multiplicar todos los precios por $\lambda = 5$ escala el valor de la dotación y el costo de las canastas por igual, dejando invariantes el conjunto presupuestario y el equilibrio real.}
\sol{99}{B}{Cada punto de la caja representa una asignación factible donde $x_{1A} + x_{1B} = X$ (ancho) y $x_{2A} + x_{2B} = Y$ (altura).}
\sol{100}{C}{Como $A$ valora más el bien $1$ en términos del bien $2$ que $B$ ($TMS^A > TMS^B$), ambos ganan si $A$ adquiere bien $1$ de $B$ y le entrega bien $2$ a una tasa intermedia entre $TMS^B$ y $TMS^A$.}
\sol{101}{A}{En el borde de la caja (por ejemplo, $O_A$, $O_B$ o con sustitutos perfectos), las curvas de indiferencia pueden no ser tangentes ($TMS^A \neq TMS^B$) y aun así ser Pareto eficientes por restricción de no negatividad.}
\sol{102}{D}{La Curva de Contrato reúne todos los Óptimos de Pareto de la caja; el Núcleo restringe ese conjunto a las asignaciones individualmente racionales ($U_A \ge U_A(W)$ y $U_B \ge U_B(W)$).}
\sol{103}{B}{Si una coalición $S$ pudiera mejorar a todos sus miembros con su dotación $\sum_{i\in S}\omega_i$, el costo a precios $p^*$ de esa canasta superaría $\sum_{i\in S} p^*\omega_i$, lo cual es imposible.}
\sol{104}{C}{Al crecer el número de agentes, las coaliciones pueden bloquear cada vez más asignaciones no competitivas; en el límite $r \to \infty$, solo sobreviven los equilibrios walrasianos.}
\sol{105}{A}{Cuando $X > Y$, el excedente $X - Y$ no aporta utilidad adicional a nadie; cualquier reparto de ese sobrante manteniendo $x_A \ge y_A$ y $x_B \ge y_B$ es Pareto eficiente (banda gruesa).}
\sol{106}{D}{La condición $TMS^A = TMS^B$ se reduce a $v_A'(x_{1A}) = v_B'(X - x_{1A})$, que determina un valor único $x_{1A}^*$ independiente de $x_{2A}$ (tramo interior vertical).}
\sol{107}{B}{La curva de oferta de $A$ une sus elecciones óptimas al rotar la recta presupuestaria alrededor de $W$. El equilibrio walrasiano ocurre en la intersección de las curvas de oferta de $A$ y $B$.}
\sol{108}{C}{Maximizar $X$ dado un nivel de $Y$ requiere tangencia entre isocuantas: $TMST_{L,K}^X = TMST_{L,K}^Y$, lo cual garantiza el mercado competitivo al igualar ambas a $w/r$.}
\sol{109}{A}{Mover recursos desde $Y$ hacia $X$ cuesta $dC = CMg_X dX + CMg_Y dY = 0$ sobre la frontera eficiente, de donde $TMT_{X,Y} = -dY/dX = CMg_X / CMg_Y$.}
\sol{110}{D}{Con intensidades factoriales distintas, reasignar factores hacia $X$ libera una proporción inadecuada de $K$ y $L$ respecto de lo que requiere $X$, generando costos de oportunidad crecientes (FPP cóncava).}
\sol{111}{B}{Si las proporciones factoriales óptimas coinciden en ambas industrias y hay RCE (o hay un solo factor), el costo de oportunidad $CMg_X/CMg_Y$ es constante y la FPP es lineal.}
\sol{112}{C}{$TMS_{X,Y} > TMT_{X,Y}$ indica que los consumidores valoran una unidad extra de $X$ en más unidades de $Y$ de las que cuesta producirla; aumentar $X$ a costa de $Y$ genera una mejora de Pareto.}
\sol{113}{A}{Al enfrentar todos los agentes los mismos precios $(p_X, p_Y, w, r)$, se concatenan $TMS^A = TMS^B = p_X/p_Y = CMg_X/CMg_Y = TMT$ y $TMST^X = TMST^Y = w/r$.}
\sol{114}{D}{La FPP entre ocio y consumo es $c = f(T - l)$, cuya pendiente en valor absoluto es $TMT_{l,c} = f'(L) = PMg_L$. En el óptimo, $TMS_{l,c} = PMg_L$ (que en descentralización iguala al salario real $w/p$).}
\sol{115}{B}{Con rendimientos crecientes, $CMg < CMe$, de modo que la regla competitiva $P = CMg$ genera pérdidas y no existe un equilibrio walrasiano precio-aceptante.}
\sol{116}{C}{Todos los consumidores enfrentan $p_X^c = (1+\tau)p_X^p$ (mantienen $TMS^A=TMS^B$) y la producción sigue sobre la FPP, pero $TMS = (1+\tau)TMT \neq TMT$.}
\sol{117}{A}{El Primer Teorema demuestra que todo equilibrio walrasiano es Pareto eficiente usando solo no saciedad local (y ausencia de externalidades/poder de mercado); la convexidad solo se exige en el Segundo Teorema.}
\sol{118}{D}{Un extremo donde un solo agente posee todos los bienes es Pareto eficiente, lo cual ilustra que eficiencia paretiana no implica justicia ni equidad distributiva.}
\sol{119}{B}{La convexidad permite separar el conjunto mejor-que del conjunto factible mediante un hiperplano de precios $p^*$; sin convexidad, el óptimo deseado puede no ser un máximo presupuestario a ningún precio.}
\sol{120}{C}{Solo las transferencias de suma alzada modifican la riqueza inicial sin introducir cuñas entre $TMS$ y $TMT$, preservando todas las condiciones de eficiencia.}
\sol{121}{A}{Al haber asimetría de información sobre la capacidad productiva individual, cualquier impuesto basado en ingresos o consumo observables deja de ser de suma alzada y crea distorsiones (trade-off eficiencia-equidad).}
\sol{122}{D}{El Primer Teorema falla cuando existen fallas de mercado (externalidades, bienes públicos, monopolio, información asimétrica o mercados incompletos), pero no por desigualdad inicial ni preferencias heterogéneas.}
\sol{123}{B}{Cuando una condición paretiana no puede cumplirse, el óptimo restringido (second-best) suele requerir desviarse deliberadamente de las demás condiciones marginales de eficiencia.}
\sol{124}{C}{En teoría, el planificador elige el punto deseado sobre la FPU (equidad) vía transferencias de suma alzada y delega al mecanismo de precios la asignación eficiente.}
\sol{125}{A}{Con mercados completos de bienes contingentes, los agentes comparten eficientemente el riesgo igualando sus tasas marginales de sustitución entre estados a los precios relativos contingentes.}
\sol{126}{D}{Construyendo una aplicación continua sobre el simplex de precios que sube precios con exceso de demanda, el Teorema de Brouwer (o Kakutani si hay correspondencias) garantiza un punto fijo $p^*$ con $z(p^*) \le 0$.}
\sol{127}{B}{SMD demuestra que la racionalidad individual no impone restricciones suficientes sobre $z(p)$ agregado (salvo continuidad, homogeneidad de grado 0 y Ley de Walras), por lo que puede haber múltiples equilibrios e inestabilidad.}
\sol{128}{C}{La propiedad de sustitubilidad bruta (gross substitutes) implica que si sube $p_j$, aumenta el exceso de demanda de los demás bienes, lo que asegura unicidad y estabilidad global del tâtonnement.}
\sol{129}{A}{El cumplimiento de WARP a nivel agregado (por ejemplo, cuando existe un consumidor representativo positivo) impide que haya múltiples precios relativos de equilibrio aislados.}
\sol{130}{D}{En economías regulares la suma de los índices de los equilibrios es $+1$ (siendo cada índice $+1$ o $-1$), lo que obliga a que exista un número finito e impar de equilibrios.}
\sol{131}{B}{Si $z_1'(p_1^*) < 0$, un precio ligeramente superior a $p_1^*$ genera exceso de oferta ($z_1 < 0$) y el subastador baja $p_1$ de vuelta hacia $p_1^*$.}
\sol{132}{C}{Cuando un agente es oferente neto de un bien, un alza de $p_1$ eleva el valor de su dotación; si ese efecto riqueza sobre su propia demanda domina al efecto sustitución, $dz_1/dp_1$ puede volverse positivo.}
\sol{133}{A}{Arrow-Debreu define una mercancía por sus características físicas, fecha, lugar y estado de la naturaleza, extendiendo el equilibrio general al análisis intertemporal y bajo incertidumbre.}
\sol{134}{D}{Si las preferencias son estrictamente monótonas, cuando $p_j \to 0$ la demanda por el bien $j$ diverge ($\|z(p)\| \to \infty$), asegurando que el equilibrio no quede en la frontera con precios nulos.}
\sol{135}{B}{Stolper-Samuelson (efecto magnificación de Jones) prueba que $\hat{w} > \hat{p}_X > \hat{p}_Y > \hat{r}$: el factor usado intensivamente en el bien que se encarece gana en términos reales de ambos bienes y el otro pierde.}
\sol{136}{C}{Por Rybczynski, a precios constantes $(w/r$ y proporciones $K/L$ fijas), absorber el trabajo extra exige expandir $X$ y retirar capital del sector $Y$, reduciendo absolutamente la producción de $Y$.}
\sol{137}{A}{Con tecnologías idénticas, RCE y sin reversión de intensidades factoriales, la igualdad de precios de bienes $(p_X, p_Y)$ determina unívocamente iguales $(w, r)$ en ambos países.}
\sol{138}{D}{La abundancia relativa de capital abarata $r/w$ en autarquía, dando ventaja comparativa en el bien intensivo en capital.}
\sol{139}{B}{Harberger mostró que gravar el capital en un sector induce sustitución hacia trabajo en $X$ (bajando $r$) y encarece $X$ contrayendo su producción (efecto producto que también presiona $r$ si $X$ es intensivo en $K$).}
\sol{140}{C}{Si $X$ es intensivo en trabajo, la contracción de $X$ (efecto producto) libera relativamente mucho trabajo y poco capital, lo que puede bajar $w/r$ y trasladar la carga al trabajo.}
\sol{141}{A}{Salvo en el caso especial de FPP lineal (costos constantes), el equilibrio general determina $p_X/p_Y$ conjuntamente por la oferta relativa (FPP) y la demanda relativa.}
\sol{142}{D}{Cada punto sobre la FPU corresponde a una asignación Pareto eficiente (Curva de Contrato y eficiencia global); ningún punto fuera de la FPU es factible.}
\sol{143}{B}{Cuando dos FPU asociadas a asignaciones distintas se cruzan, $B$ puede dominar a $A$ según Kaldor-Hicks y simultáneamente $A$ dominar a $B$, generando intransitividad.}
\sol{144}{C}{Para $\rho = 1$ es Bentham (cero aversión a desigualdad en utilidades); para $\rho \to 0$ es Nash ($U_1 U_2$); para $\rho \to -\infty$ es Rawls ($\min\{U_1, U_2\}$, aversión infinita).}
\sol{145}{A}{Arrow probó que cualquier regla transitiva y completa que cumpla dominio universal, Pareto e IIA es necesariamente dictatorial.}
\sol{146}{D}{IIA prohíbe que la presencia o posición relativa de una tercera alternativa $z$ invierta la clasificación social entre $x$ e $y$ (reglas como el conteo de Borda violan IIA).}
\sol{147}{B}{El Teorema de Black (votante mediano) muestra que si las preferencias son unimodales en un espacio unidimensional (y hay número impar de votantes), el pico del votante mediano vence a cualquier otra alternativa.}
\sol{148}{C}{Maximizar $W(U_A, U_B)$ sujeto a la FPU requiere tangencia (en solución interior): el cociente de ponderaciones marginales $(\partial W/\partial U_A)/(\partial W/\partial U_B)$ iguala la pendiente de la FPU.}
\sol{149}{A}{El Segundo Teorema sugiere mantener las señales de precio eficientes ($P = CMg$) y abordar la equidad vía transferencias directas de suma alzada (bonos focalizados), evitando distorsiones por subsidios cruzados excesivos.}
\sol{150}{D}{El Análisis Costo-Beneficio estándar aprueba proyectos con beneficio neto positivo ($\Delta W > 0$), lo que corresponde a una mejora potencial de Pareto (criterio de Kaldor-Hicks).}
\sol{151}{B}{La función de bienestar de Rawls ($W = \min_i U_i$) asigna prioridad absoluta a mejorar la situación de los individuos con menor nivel de bienestar, como los hogares sin acceso básico al agua.}
\sol{152}{C}{Sin mercados completos de bienes contingentes (Arrow-Debreu) para estados secos y húmedos, los agentes no pueden transferir eficientemente el riesgo hidrológico.}
\sol{153}{A}{Las economías de escala en redes (monopolio natural) y las externalidades positivas del saneamiento (salud pública) o negativas de vertimientos violan los supuestos del Primer Teorema.}
\sol{154}{A}{Por la Ley de Walras, $p_1 z_1 + p_2 z_2 + p_3 z_3 = 0 \implies 2(10) + 3(-5) + 5z_3 = 0 \implies 5 + 5z_3 = 0 \implies z_3 = -1$.}
\sol{155}{B}{Igualando $z_1 = 0$: $6(1) - 2p_1^* = 0 \implies p_1^* = 3$. Al reemplazar $(\lambda p_1, \lambda p_2)$, el factor $\lambda$ se cancela, confirmando homogeneidad de grado $0$.}
\sol{156}{C}{Igualando $TMS_A = TMS_B$: $\frac{1}{2}\frac{y_A}{x_A} = 2\frac{12 - y_A}{12 - x_A}$. Con $x_A = 4$ ($x_B = 8$): $\frac{y_A}{8} = \frac{12 - y_A}{4} \implies y_A = 24 - 2y_A \implies y_A = 8$.}
\sol{157}{D}{Con $p_y=1$, $I_A = 12p_x$ y $I_B = 6$, dando $y_A = 6p_x$ y $y_B = 3$. Vaciando el mercado de $y$: $6p_x + 3 = 6 \implies p_x^* = 0.5$, $x_A^* = 6$ y $y_A^* = 3$.}
\sol{158}{A}{Con $p_y=1$, $x_A = \frac{1}{4}\frac{4p_x+4}{p_x} = 1 + \frac{1}{p_x}$ y $x_B = \frac{3}{4}\frac{4p_x+12}{p_x} = 3 + \frac{9}{p_x}$. Igualando $x_A + x_B = 8$: $4 + \frac{10}{p_x} = 8 \implies p_x = 2.50$.}
\sol{159}{B}{De $TMS_i = p_x$ se obtiene $x_A = 6/p_x$ y $x_B = 10/p_x$. Vaciando el mercado de $x$: $16/p_x = 8 \implies p_x^* = 2$, resultando $x_A^* = 3$ y $x_B^* = 5$ (con $y_A^*=8, y_B^*=12$).}
\sol{160}{C}{En la dotación inicial $U_A(W) = U_B(W) = 9$. Sobre la curva de contrato ($y_A = x_A$), la racionalidad individual exige $x_A^2 \ge 9 \implies x_A \ge 3$ y $(10-x_A)^2 \ge 9 \implies x_A \le 7$.}
\sol{161}{D}{En $A=(7,7)$, $TMS_A = TMS_B = 1$, luego $p_x^* = p_y = 1$. A esos precios, el ingreso inicial de $A$ es $4+2=6$ y su canasta objetivo cuesta $7+7=14$, requiriendo una transferencia $T = 14 - 6 = 8$.}
\sol{162}{A}{Como la dotación total de $x$ ($10$) excede a la de $y$ ($6$) y ambos requieren proporción $1:1$, el bien $x$ es abundante (bien libre), luego $p_x^* = 0$. Así $I_A = 4$, $y_A^* = 4$ y $U_A^* = 4$.}
\sol{163}{B}{La eficiencia técnica exige $TMST_X = TMST_Y$: $\frac{K_X}{L_X} = \frac{1}{2}\frac{90 - K_X}{90 - L_X}$. Con $L_X = 45$ ($L_Y = 45$): $K_X = \frac{1}{2}(90 - K_X) \implies 1.5K_X = 45 \implies K_X = 30$.}
\sol{164}{C}{Diferenciando la FPP: $2X dX + 8Y dY = 0 \implies TMT = X/(4Y)$. Igualando a $TMS = Y/X$ resulta $X^2 = 4Y^2$. Sustituyendo en la FPP: $8Y^2 = 200 \implies Y^*=5$, $X^*=10$ y $p_X/p_Y = 5/10 = 0.5$.}
\sol{165}{D}{La FPP es $X/2 + Y/4 = 30 \iff 2X + Y = 120$, por lo que $TMT = p_X/p_Y = 2$. Igualando $TMS = Y/X = 2$ se tiene $Y = 2X$, dando $4X = 120 \implies X^* = 30$ y $Y^* = 60$.}
\sol{166}{A}{En $(4, 8)$: $TMT = CMg_X/CMg_Y = 8/32 = 0.25$ y $TMS = 8/4 = 2.0$. Como $TMS > TMT$, la valoración marginal de $X$ supera su costo de oportunidad: conviene producir más $X$ y menos $Y$.}
\sol{167}{B}{Igualando $TMS_{l,c} = PMg_L$: $\frac{c}{l} = \frac{3}{\sqrt{L}}$. Sustituyendo $c = 6\sqrt{L}$ y $l = 27 - L$: $\frac{6\sqrt{L}}{27 - L} = \frac{3}{\sqrt{L}} \implies 2L = 27 - L \implies L^* = 9$ y $c^* = 6(3) = 18$.}
\sol{168}{C}{Sobre la FPP, $TMT = p_X^p/p_Y = 1/2$. Con impuesto del $100\%$, $p_X^c/p_Y = 1$, por lo que $TMS = Y/X = 1 \implies Y = X$. Sustituyendo en $X + 2Y = 120$ resulta $X_t = 40$ e $Y_t = 40$.}
\sol{169}{D}{Resolviendo $2w_1 + r_1 = 13$ y $w_1 + 2r_1 = 11$: multiplicando la primera por $2$ y restando la segunda da $3w_1 = 15 \implies w_1 = 5$ y $r_1 = 3$ (frente a $w_0=3, r_0=4$).}
\sol{170}{A}{Restando las ecuaciones en variaciones: $2\Delta X + \Delta Y = 6$ y $\Delta X + 2\Delta Y = 0 \implies \Delta X = -2\Delta Y \implies -3\Delta Y = 6 \implies \Delta Y = -2$ y $\Delta X = +4$.}
\sol{171}{B}{La tangencia entre la recta de isobienestar y la FPU exige $\frac{3}{1} = \frac{U_A}{U_B} \implies U_A = 3U_B$. Sustituyendo en la FPU: $10U_B^2 = 160 \implies U_B^* = 4$ y $U_A^* = 12$.}
\sol{172}{C}{Igualando la pendiente de $W$ ($U_B/U_A$) a la pendiente de la FPU ($2$): $U_B = 2U_A$. Sustituyendo en $2U_A + U_B = 120$ da $4U_A = 120 \implies U_A^* = 30$ y $U_B^* = 60$.}
\sol{173}{D}{El óptimo rawlsiano se alcanza en el vértice $U_A = U_B$ sobre la FPU: $U_A + 3U_A = 120 \implies 4U_A = 120 \implies U_A^* = U_B^* = 30$.}
\sol{174}{A}{Igualando la tasa marginal de sustitución social $(U_B/U_A)^{1/2}$ a la pendiente de la FPU ($1/4$): $U_B/U_A = 1/16 \implies U_A = 16U_B$. En la FPU: $20U_B = 100 \implies U_B^* = 5$ y $U_A^* = 80$.}
\sol{175}{B}{Para que ambos grupos queden estrictamente mejor, $B$ debe recibir más que su pérdida ($T > 110$) y $A$ debe pagar menos que su ganancia ($T < 180$), es decir, $110 < T < 180$.}
\sol{176}{C}{Por equimarginalidad, $100 - 2q_1 = 85 - (60 - q_1) \implies 100 - 2q_1 = 25 + q_1 \implies 3q_1 = 75 \implies q_1^* = 25$, $q_2^* = 35$ y $\lambda^* = 100 - 50 = 50$.}
\sol{177}{D}{En $(45, 15)$, $BMg_1 = 10$ y $BMg_2 = 70$. Al transferir $20$ unidades de $1$ a $2$ hasta igualar ambos en $50$, el triángulo de ganancia neta es $\Delta W = \frac{1}{2}(70 - 10)(20) = 600$.}
\sol{178}{A}{Por el Teorema del Votante Mediano (preferencias unimodales con número impar de votantes), el ganador de Condorcet es la mediana de los picos ($S/\ 50$), no el promedio ($S/\ 61$).}
\sol{179}{B}{Como ambos tienen preferencias idénticas homotéticas, en equilibrio $c_{iH}/c_{iS} = p_S/p_H = C_H/C_S = (12+4)/(4+4) = 16/8 = 2.0$ (el agua es dos veces más cara en sequía).}
\sol{180}{C}{Con $p_x'=4$, el hogar elige $x'=9, y'=36$ ($U'=18$) y el subsidio cuesta $(9-4)(9) = 45$. A precios $(9,1)$, el gasto mínimo para $U'=18$ es $e = 2\sqrt{9\cdot 1}(18) = 108$, requiriendo una suma alzada de $108 - 72 = 36$ (ahorro de $9$).}


\end{document}
